% Auto-generated by Mathematica/CubeGalerkin27.wl with hand-edited
% \begin{aligned}[t]...\end{aligned} line breaks in Tables 2--3 to prevent
% long closed-form Mp expressions from overflowing the page. If you
% regenerate this file, reapply the aligned[t] wrappers to the overflowing
% rows (see git log for this file).
\subsection{Verified analytical master integrals}
\label{sec:results}

The verification at degree 0 was reported in Section~\ref{sec:tier1}.
We extend the table to include all parity-allowed master integrals up to degree~6:

\begin{table}[h]
\centering
\begin{tabular}{lll}
\toprule
$(p,q,r)$ & $\mathcal{M}_{p,q,r} = \int_{[-1,1]^3}\!\frac{x^p y^q z^r}{|\bm{r}|}\,dV$ & numerical \\
\midrule
$(0,0,0)$ & $-2\pi-\tfrac{4}{3}\log(70226-40545\sqrt{3})$ & $9.5203095$ \\
$(2,0,0)$ & $\tfrac{1}{9}\bigl(6\sqrt{3}-\pi+\log(3650401+2107560\sqrt{3})\bigr)$ & $2.5615786$ \\
$(2,2,0)$ & $\begin{aligned}[t]&\tfrac{1}{135}\bigl(108\sqrt{3}+18\pi+63\log 2+20\log(26-15\sqrt{3})\\&\quad{}-60\log(\sqrt{3}-1)-66\log(1+\sqrt{3})-90\coth^{-1}\sqrt{3}\bigr)\end{aligned}$ & $0.75095335$ \\
$(4,0,0)$ & $\begin{aligned}[t]&-\tfrac{2}{45}\bigl(12\sqrt{3}+7\pi-42\log 2\\&\quad{}+99\log(\sqrt{3}-1)-15\log(1+\sqrt{3})\bigr)\end{aligned}$ & $1.4351485$ \\
$(2,2,2)$ & $\begin{aligned}[t]&\tfrac{1}{1260}\bigl(1008\sqrt{3}-42\pi+522\log 2+441\log(\sqrt{3}-1)\\&\quad{}-1395\log(1+\sqrt{3})-84\log(2+\sqrt{3})-90\log(11\sqrt{3}-19)\\&\quad{}+35\log(26+15\sqrt{3})-672\coth^{-1}\sqrt{3}\bigr)\end{aligned}$ & $0.22735457$ \\
$(4,2,0)$ & $\begin{aligned}[t]&\tfrac{1}{630}\bigl(-48\sqrt{3}+17\pi\\&\quad{}+12\tanh^{-1}\tfrac{21252634831\sqrt{3}}{36810643322}\bigr)\end{aligned}$ & $0.42942069$ \\
$(6,0,0)$ & $\tfrac{1}{42}\bigl(-5\pi-4\log 2+8(3\sqrt{3}+\log(5+3\sqrt{3}))\bigr)$ & $0.99201788$ \\
\bottomrule
\end{tabular}
\caption{Closed-form values of the scalar master integrals $\mathcal{M}_{p,q,r}$ over the unit cube.  Cubic symmetry equates entries that are permutations of each other; only canonical orderings $p\ge q \ge r$ are listed.}
\label{tab:tier1-masters}
\end{table}

The same family of integrals over the positive octant $[0,1]^3$ supplies the body-bilinear-form building blocks; we collect them in Table~\ref{tab:tier2-masters}.

\begin{table}[h]
\centering
\begin{tabular}{lll}
\toprule
$(p,q,r)$ & $\mathcal{M}^+_{p,q,r} = \int_{[0,1]^3}\!\frac{x^p y^q z^r}{|\bm{r}|}\,dV$ & numerical \\
\midrule
$(0,0,0)$ & $-\frac{\pi }{4}-\frac{1}{6} \log \left(70226-40545 \sqrt{3}\right)$ & $1.1900387$ \\
$(1,0,0)$ & $\begin{aligned}[t]&\tfrac{1}{36}\bigl(-2\pi-12\bigl(\sqrt{2}-\sqrt{3}+\sinh^{-1}(1)\bigr)\\&\quad{}+\cosh^{-1}(189750626)+18\coth^{-1}\sqrt{3}\bigr)\end{aligned}$ & $0.51559356$ \\
$(1,1,0)$ & $\begin{aligned}[t]&\tfrac{1}{12}\bigl(1-7\sqrt{2}+6\sqrt{3}+\log(5\sqrt{2}-7)\\&\quad{}-3\log(2-\sqrt{3})\bigr)\end{aligned}$ & $0.23329690$ \\
$(2,0,0)$ & $\begin{aligned}[t]&\tfrac{1}{144}\bigl(12\sqrt{3}-2\pi\\&\quad{}+\log(26650854921601+15386878263120\sqrt{3})\bigr)\end{aligned}$ & $0.32019732$ \\
$(1,1,1)$ & $\tfrac{1}{5}\bigl(1-4\sqrt{2}+3\sqrt{3}\bigr)$ & $0.10785963$ \\
$(2,1,0)$ & $\begin{aligned}[t]&\tfrac{1}{2880}\bigl(-504\sqrt{2}+672\sqrt{3}+32\pi+72\log(\sqrt{3}-1)\\&\quad{}-252\log(1+\sqrt{3})+90\log(2(2+\sqrt{3}))-261\sinh^{-1}(1)\\&\quad{}+40\cosh^{-1}(26)+45\coth^{-1}\sqrt{2}\bigr)\end{aligned}$ & $0.14741068$ \\
\bottomrule
\end{tabular}
\caption{Closed-form values of the positive-octant master integrals $\mathcal{M}^+_{p,q,r}$ over $[0,1]^3$. Up to canonical ordering $p\ge q \ge r$ and $p+q+r\le 3$ they exhaust all polynomial weights needed by the Path-B body bilinear form on the constant-and-linear subspace.}
\label{tab:tier2-masters}
\end{table}

\subsection{Body bilinear-form building blocks}
\label{sec:results-tier3}

For the constant subspace of $T_{27}$ the diagonal and off-diagonal entries of the body bilinear form $\mathcal{B}_{\mathrm{body}}$ reduce, after the centre-of-mass / relative-coordinate substitution and an octant-symmetry restriction, to three positive scalars.
Define
\begin{align}
  i_A      &:= \int_{[0,1]^3}\!(1-u_1)(1-u_2)(1-u_3)\,\frac{1}{|\bm{u}|}\,d\bm{u},\\
  i_{B,11} &:= \int_{[0,1]^3}\!(1-u_1)(1-u_2)(1-u_3)\,\frac{u_1^2}{|\bm{u}|^3}\,d\bm{u},\\
  i_{B,12} &:= \int_{[0,1]^3}\!(1-u_1)(1-u_2)(1-u_3)\,\frac{u_1 u_2}{|\bm{u}|^3}\,d\bm{u}.
\end{align}
The first reduces by linearity to a signed combination of $\mathcal{M}^+_{p,q,r}$ values (Table~\ref{tab:tier2-masters}):
\begin{equation}
  i_A \;=\; \mathcal{M}^+_{0,0,0} \;-\; 3\mathcal{M}^+_{1,0,0} \;+\; 3\mathcal{M}^+_{1,1,0} \;-\; \mathcal{M}^+_{1,1,1}
       \;\approx\; 0.2352890805487075.
  \label{eq:iA-closed}
\end{equation}

The remaining two are obtained by integrating once by parts in the $u_1$ direction; the singular kernel $|\bm{u}|^{-3}$ collapses to a $|\bm{u}|^{-1}$-class kernel and a clean two-dimensional residual:
\begin{align}
  i_{B,11} &\;=\; i_A \;-\; J_{11},\qquad
  J_{11} \;=\; \int_{[0,1]^2}\!(1-u_2)(1-u_3)\Bigl[\sqrt{1+u_2^2+u_3^2} - \sqrt{u_2^2+u_3^2}\,\Bigr]\,du_2\,du_3,\\
  i_{B,12} &\;=\; \int_{[0,1]^2}\!(1-u_2)(1-u_3)\,u_2\,\Bigl[\frac{1}{\sqrt{u_2^2+u_3^2}} \;-\; \operatorname{arcsinh}\!\Bigl(\frac{1}{\sqrt{u_2^2+u_3^2}}\Bigr)\Bigr]\,du_2\,du_3.
\end{align}

Mathematica evaluates both two-dimensional integrals in closed form:
\begin{equation*}
\resizebox{\linewidth}{!}{$\displaystyle
J_{11} \;=\; \frac{1}{240} \left(8+8 \sqrt{2}-16 \sqrt{3}-55 \log (2)+15 \log \left(3+2 \sqrt{2}\right)+110 \log \left(1+\sqrt{3}\right)-15 \log \left(2+\sqrt{3}\right)-80 \tan ^{-1}\left(\frac{3}{8}\right)+80 \tan ^{-1}\left(\frac{1}{183} \left(96-73 \sqrt{3}\right)\right)+10 \sinh ^{-1}(1)\right)
$}
\end{equation*}
with $J_{11} \approx 0.1568593870324717$, and
\begin{equation*}
\resizebox{\linewidth}{!}{$\displaystyle
i_{B,12} \;=\; \frac{1}{240} \left(-8 \sqrt{2}+4 \sqrt{3}+20 \log (2)+5 \log (256)+40 \log \left(1+\sqrt{2}\right)+22 \log \left(3+2 \sqrt{2}\right)-120 \log \left(1+\sqrt{3}\right)-22 \log \left(2+\sqrt{3}\right)+20 \tan ^{-1}\left(\frac{1}{2}\right)+16 \tan ^{-1}\left(\frac{5}{8}\right)-40 \tan ^{-1}\left(\frac{3}{2}\right)+40 \tan ^{-1}\left(8+\frac{13}{\sqrt{3}}\right)-20 \tan ^{-1}\left(\frac{1}{11} \left(8-5 \sqrt{3}\right)\right)-16 \tan ^{-1}\left(\frac{1}{167} \left(160+89 \sqrt{3}\right)\right)+4 \sinh ^{-1}(1)+36 \sinh ^{-1}\left(\frac{1}{\sqrt{2}}\right)\right)
$}
\end{equation*}
with $i_{B,12} \approx 0.04970171296728090$.

A non-trivial consistency check is supplied by the trace identity
\begin{equation*}
  3\,i_{B,11} \;=\; \int_{[0,1]^3}\!(1-u_1)(1-u_2)(1-u_3)\,\frac{u_1^2+u_2^2+u_3^2}{|\bm{u}|^3}\,d\bm{u} \;=\; i_A,
\end{equation*}
which is verified analytically and numerically to working precision.


\subsection{The 9{\times}9 body bilinear matrix on $T_9$}
\label{sec:results-body9}

With the master integrals of Sections~\ref{sec:results} and the three scalar building blocks $i_A$, $i_{B,11}$, $i_{B,12}$, the body bilinear form $\mathcal{B}_9 : T_9\times T_9 \to \mathbb{R}$ is assembled entry-by-entry from
\begin{equation}
  \mathcal{B}_9[i,j] \;=\; 32\,\sum_{p,q=1}^{3}\Bigl\{A\,\delta_{pq}\,K_{1}[\Pi_{ij,pq}] \,+\, B\,K_{3}[\Pi_{ij,pq};\,p,q]\Bigr\},
\end{equation}
where $\Pi_{ij,pq}(\bm{u})$ is the residual polynomial obtained by (a) substituting $\phi^{(i)}_p(\bm{r})\phi^{(j)}_q(\bm{r}')$, (b) integrating the centre-of-mass variable $\bm{\xi}$ over the overlap box, and (c) dividing out the standard tent factor $(1-u_1)(1-u_2)(1-u_3)$.  The prefactor $32 = 4\cdot 8/1$ absorbs the $s\mapsto 2u$ Jacobian $8$, the Green-tensor rescaling $G(2\bm{u})=\tfrac12 [A/|\bm{u}| + B\,u_p u_q/|\bm{u}|^3]$, and the symmetric-box factor $2/1$.  A further factor of $8$ appears on parity-even entries when we restrict to the positive octant $[0,1]^3$, so a diagonal displacement-block entry is proportional to $256$.

Under cubic symmetry the $T_9$-basis decomposes into four irreducible blocks, each of which is characterised by at most two independent scalars:
\begin{itemize}
  \item \textbf{Displacement block} ($i,j\in\{1,2,3\}$): diagonal, with a single scalar.
  \item \textbf{Displacement--strain block} ($i\in\{1,2,3\}$, $j\in\{4,\ldots,9\}$): identically zero by parity.
  \item \textbf{Axial-strain block} ($i,j\in\{4,5,6\}$): diagonal entry $\alpha$, off-diagonal entry $\beta$.
  \item \textbf{Axial--shear block} ($i\in\{4,5,6\}$, $j\in\{7,8,9\}$): identically zero by parity.
  \item \textbf{Shear-strain block} ($i,j\in\{7,8,9\}$): diagonal entry $\gamma$, off-diagonals identically zero by parity.
\end{itemize}

\subsubsection*{Displacement block}
All three displacement-block diagonal entries are equal and reduce in closed form to
\begin{equation}
\resizebox{\linewidth}{!}{$\displaystyle
  \mathcal{B}_9[1,1] \;=\; \mathcal{B}_9[2,2] \;=\; \mathcal{B}_9[3,3]
     \;=\; 256\,\bigl(A\,i_A \;+\; B\,i_{B,11}\bigr)
     \;=\; 60.23400462046913\,A \;+\; 20.07800154015638\,B,
$}
  \label{eq:B9disp}
\end{equation}

with off-diagonal entries vanishing by parity.

\subsubsection*{Axial-strain block}
The diagonal entry $\alpha := \mathcal{B}_9[4,4] = \mathcal{B}_9[5,5] = \mathcal{B}_9[6,6]$ splits into
\begin{equation}
\resizebox{\linewidth}{!}{$\displaystyle
  \alpha \;=\; \tfrac{256}{3}\Bigl[A\bigl(i_A - 2 K_1^{(2,0,0)}\bigr)
     \;+\; B\bigl(i_{B,11} - 2 K_3^{(2,0,0)}\bigr)\Bigr]
     \;=\; 15.37252824403423\,A \;+\; 3.949761457945144\,B,
$}
  \label{eq:alpha-axial}
\end{equation}
where
\begin{equation*}
\resizebox{\linewidth}{!}{$\displaystyle
  K_1^{(2,0,0)} \;:=\; \int_{[0,1]^3}(1-u_1)(1-u_2)(1-u_3)\,\frac{u_1^2}{|\bm{u}|}\,d\bm{u},
  \quad K_3^{(2,0,0)} \;:=\; \int_{[0,1]^3}(1-u_1)(1-u_2)(1-u_3)\,\frac{u_1^4}{|\bm{u}|^3}\,d\bm{u}.
$}
\end{equation*}
The off-diagonal entry $\beta := \mathcal{B}_9[4,5]$ is purely a $B$-channel correction,
\begin{equation}
  \beta \;=\; -256\,B\,K_{3,\text{off}}^{(1,1,0)}
     \;=\; 0\,A \;-\; 1.471925680511371\,B,
  \label{eq:beta-axial}
\end{equation}
where $K_{3,\text{off}}^{(1,1,0)} = \int_{[0,1]^3}(1-u_1)(1-u_2)(1-u_3)\,u_1 u_2\,\tfrac{u_1 u_2}{|\bm{u}|^3}\,d\bm{u}$.  Note that $\beta$ vanishes in the $A$-channel because the integrand of the $A$ piece is odd in $(u_1,u_2)$.

\subsubsection*{Shear-strain block}
The diagonal entry $\gamma := \mathcal{B}_9[7,7] = \mathcal{B}_9[8,8] = \mathcal{B}_9[9,9]$ evaluates to
\begin{equation}
\resizebox{\linewidth}{!}{$\displaystyle
  \gamma \;=\; \tfrac{128}{3}\Bigl[A\bigl(i_A - 2 K_1^{(2,0,0)}\bigr) + B\bigl(i_{B,11} - 2 K_3^{(0,2,0)}\bigr)\Bigr]
     \;-\; 128\,B\,K_{3,\text{off}}^{(1,1,0)}
     \;=\; 7.686264122017116\,A \;+\; 2.119728856266587\,B,
$}
  \label{eq:gamma-shear}
\end{equation}
where $K_3^{(0,2,0)} = \int_{[0,1]^3}(1-u_1)(1-u_2)(1-u_3)\,\tfrac{u_2^2 u_1^2}{|\bm{u}|^3}\,d\bm{u}$.  The relation $\gamma^{(A)} = \tfrac12\,\alpha^{(A)}$ is an exact identity: the $A$-channel of $\gamma$ is exactly half the $A$-channel of $\alpha$, coming from the bilinearity of the two-component shear basis $\phi^{(7)} = (0,r_3/2,r_2/2)$.

Numerical verification of cubic symmetry (all six independent identities) was performed to $10^{-20}$ precision; see the Mathematica script output.


\subsection{The 18{\times}18 quad--quad block on \texorpdfstring{$T_{27}$}{T27}}
\label{sec:results-body27quad}

Extending $T_9$ to $T_{27}$ adjoins the 18 quadratic vector fields
\begin{equation*}
  \phi^{(9 + 6(k-1) + m)}(\br) \;=\; q_m(\br)\,\be_k,
  \qquad k \in \{1,2,3\},\;\; m \in \{1,\ldots,6\},
\end{equation*}
where the six quadratic monomials in canonical Voigt order are
\begin{equation*}
  q_1 = r_1^2,\quad q_2 = r_2^2,\quad q_3 = r_3^2,\quad
  q_4 = r_2 r_3,\quad q_5 = r_1 r_3,\quad q_6 = r_1 r_2.
\end{equation*}
We refer to $q_{1..3}$ as the squared ($S$) modes and to $q_{4..6}$ as the cross ($X$) modes.  On the $18\times 18$ quad--quad sub-block, the body bilinear form decomposes into an $A$-channel and a $B$-channel,
\begin{equation}
  \mathcal{B}_{27}^{\text{quad}}[i,j]
  \;=\; 32\,\bigl(a_{ij}\,A \;+\; b_{ij}\,B\bigr),
  \label{eq:B27quad-split}
\end{equation}
where the scalars $a_{ij}, b_{ij}$ depend only on the orbit of $(i,j)$ under the combined O$_h$ action on the three coordinate axes and the $i\leftrightarrow j$ symmetry.

Enumerating all $18^2 = 324$ index pairs and partitioning by this action yields exactly \textbf{34 orbits}.  Of these, $\textbf{20 orbits}$ (total multiplicity $222$) are identically zero in both channels: each such vanishing is a parity cancellation forced by cubic symmetry (the integrand is odd in at least one coordinate axis under the $[-1,1]^3$ integration).  The remaining \textbf{14 orbits} (total multiplicity $102$) carry nontrivial values, and are listed in Table~\ref{tab:tier6-orbits}.

\begin{table}[h]
\centering\small
\resizebox{\linewidth}{!}{%
\begin{tabular}{llrrr}
\toprule
Representative $(i,j)$ & Structure & mult & $A$-value & $B$-value \\
\midrule
$(10,10)$ & $S$-$S$ diag, axis-matched   & $3$  & $17.126396135464$ & $ 7.086377754857$ \\
$(11,11)$ & $S$-$S$ diag, axis-unmatched & $6$  & $17.126396135464$ & $ 5.020009190303$ \\
$(10,11)$ & $S$-$S$ same-block off-diag, one matched & $12$ & $14.698679523965$ & $ 5.444954908314$ \\
$(11,12)$ & $S$-$S$ same-block off-diag, both unmatched & $6$  & $14.698679523965$ & $ 3.808769707336$ \\
$(13,13)$ & $X$-$X$ diag, no-factor-matched & $3$  & $ 3.988696969908$ & $ 1.644075434840$ \\
$(14,14)$ & $X$-$X$ diag, one-factor-matched & $6$  & $ 3.988696969908$ & $ 1.172310767534$ \\
\midrule
$(10,17)$ & $S$-$S$ cross-block (type Ia) & $6$  & $0$               & $ 0.654189191338$ \\
$(11,16)$ & $S$-$S$ cross-block (type Ib) & $6$  & $0$               & $ 0.654189191338$ \\
$(14,26)$ & $X$-$X$ cross-block (type I)  & $6$  & $0$               & $ 0.654189191338$ \\
$(13,20)$ & $X$-$X$ cross-block (type II) & $6$  & $0$               & $-0.389578544263$ \\
$(14,19)$ & $X$-$X$ cross-block (type III)& $6$  & $0$               & $-0.389578544263$ \\
$(10,21)$ & $S$-$X$ cross-block (type II) & $12$ & $0$               & $-0.216414577671$ \\
$(11,21)$ & $S$-$X$ cross-block (type III)& $12$ & $0$               & $-0.216414577671$ \\
\bottomrule
\end{tabular}%
}
\caption{The 13 non-vanishing orbits of the $18{\times}18$ quad--quad block $\mathcal{B}_{27}^{\text{quad}}$.  Indices $i,j$ are global ($10\le i,j\le 27$).  The $A$- and $B$-values in each row are the two numerical coefficients of $A$ and $B$ in \eqref{eq:B27quad-split}.  Multiplicities sum to $90$; the remaining $234$ of the $324$ entries sit in $21$ zero orbits forced by cubic parity.  Values are exact to all $12$ digits displayed; the underlying high-precision data (\texttt{WorkingPrecision}$\to 30$, see Section~\ref{sec:results-high-precision}) carries $\sim 28$ correct digits.}
\label{tab:tier6-orbits}
\end{table}

\paragraph{Independent scalars.}
The 13 non-zero orbits collapse to \textbf{3 distinct $A$-scalars and 9 distinct $B$-scalars}:
\begin{align*}
  \alpha_1 &\approx 17.126396135464, &
  \alpha_2 &\approx 14.698679523965, &
  \alpha_3 &\approx  3.988696969908, \\[4pt]
  \beta_1 &\approx  7.086377754857, &
  \beta_2 &\approx  5.020009190303, &
  \beta_3 &\approx  5.444954908314, \\
  \beta_4 &\approx  3.808769707336, &
  \beta_5 &\approx  1.644075434840, &
  \beta_6 &\approx  1.172310767534, \\
  \beta_8 &\approx  0.654189191338, &
  \beta_9 &\approx -0.389578544263, &
  \beta_{10} &\approx -0.216414577671.
\end{align*}
The numbering omits $\beta_7$ deliberately: the orbit $(11,17)$ that previously appeared to give $\beta_7 \approx 0.692769$ is, at \texttt{WorkingPrecision}$\to 30$, identically zero in both channels (see paragraph below and Section~\ref{sec:results-high-precision}).  Across all $324$ entries the observed within-orbit spread is $0$ exactly in the $A$-channel and $< 10^{-27}$ in the $B$-channel, so these 12 scalars completely determine the $18\times 18$ block.  All values are exported to \texttt{Mathematica/CubeT6ScalarValues\_HighPrec.wl}.

\paragraph{Symmetry structure and accidental identities.}
The $A$-channel vanishes on every cross-direction orbit ($k_i\neq k_j$) and on every $S$-$X$ orbit: this reflects the fact that the $A$ part of the elastic Green's tensor is proportional to $\delta_{ij}/|\mathbf{r}-\mathbf{r}'|$ (an isotropic trace piece), so it only connects modes carrying identical Cartesian vector components and with an even polynomial overlap.  The $B$-channel, which carries the $(\mathbf{r}{-}\mathbf{r}')_i(\mathbf{r}{-}\mathbf{r}')_j/|\mathbf{r}{-}\mathbf{r}'|^3$ anisotropic piece, is nonzero on all $14$ orbits.

An additional structural observation from Table~\ref{tab:tier6-orbits} is that the $7$ surviving cross-direction orbits collapse to only \emph{three} distinct $B$-values: $\beta_8=0.6542$ (appearing in three orbits), $\beta_9=-0.3896$ (appearing in two orbits), and $\beta_{10}=-0.2164$ (appearing in two orbits).  These collapses are identities to $\ge 27$-digit precision (\texttt{WorkingPrecision}$\to 30$) that are not a direct consequence of the O$_h$-orbit decomposition alone: our canonical labelling distinguishes $(10,17)$, $(11,16)$, and $(14,26)$ as separate O$_h$ orbits, yet they integrate to the same number.  Together with the unconditional vanishing of orbit $(11,17)$ (described next), they reflect a deeper structural identity of the cube Green-tensor double integral; closed-form verification via PSLQ on the high-precision values of Section~\ref{sec:results-high-precision} is the natural next step.

\paragraph{The vanishing of orbit $(11,17)$.}
A separate, qualitative finding of the high-precision recomputation is that the $S$-$X$ cross-block orbit at $(11,17)$ ($S$-$X$ type I, multiplicity $12$) is \emph{exactly zero in both channels}.  Earlier \texttt{MachinePrecision} assemblies returned $B \approx 0.692769$ for this orbit, which we previously labelled $\beta_7$; the overnight \texttt{WorkingPrecision}$\to 30$ recomputation collapses this value to $|B| < 10^{-27}$, consistent with an exact algebraic cancellation rather than residual noise.  We have therefore reclassified $(11,17)$ as a zero orbit and dropped $\beta_7$ from the list of independent scalars.  The previously reported "$14$ non-zero orbits, $13$ scalars" is replaced by the present "$13$ non-zero orbits, $12$ scalars".  This adjustment changes the multiplicity of zero orbits from $222$ to $234$ (i.e.\ $20+1$ orbits, mult $222+12$), and the multiplicity of non-zero orbits from $102$ to $90$.

\paragraph{Extraction method and precision.}
{\sloppy
The scalars $\alpha_i, \beta_j$ were initially computed by \texttt{Mathematica/\allowbreak CubeT6Block.wl}, which assembles $\mathcal{B}_{27}^{\text{quad}}$ via \texttt{NIntegrate} at \texttt{MachinePrecision} with \texttt{PrecisionGoal}\allowbreak$\to 12$.  That configuration completes all $324$ entries in $\sim 150$ s and cross-checks against the $T_9$ Tier-5 displacement diagonal \eqref{eq:B9disp} to better than $3\times 10^{-12}$, but it is just at the edge of double precision: the trailing $1$--$3$ decimal digits of the $\alpha_i,\beta_j$ are noise, and orbit $(11,17)$ in particular is misreported as having $B\approx 0.6928$ rather than its true value $0$.\par

The dedicated overnight high-precision recomputation \texttt{Mathematica/\allowbreak CubeT6OvernightHighPrec.wl} ran the same $14$ orbit representatives at \texttt{WorkingPrecision}\allowbreak$\to 30$ on $2026$-$04$-$08$, completing in $\sim 1217$ s of total \texttt{NIntegrate} time.  Its results are exported to \texttt{Mathematica/\allowbreak CubeT6ScalarValues\_\allowbreak HighPrec.wl} and tabulated in Section~\ref{sec:results-high-precision}.  The high-precision pass:\par}
\begin{enumerate}[label=(P\arabic*),leftmargin=2em]
\item confirms the $12$-digit values of the surviving $13$ orbits (corrections at the $11$th--$12$th decimal where machine-precision drift was visible);
\item shrinks the within-orbit spread of the $\alpha$- and $\beta$-channels to $< 10^{-27}$, raising the cubic-symmetry consistency check from $\sim 10^{-12}$ to $\sim 10^{-27}$;
\item promotes orbit $(11,17)$ from "$\beta_7 \approx 0.6928$" to a confirmed structural zero;
\item supplies the $\ge 28$ correct digits needed for PSLQ-based closed-form reconstruction of the $\alpha_i$ and $\beta_j$ in the same spirit as \eqref{eq:alpha-axial}--\eqref{eq:gamma-shear}.
\end{enumerate}
The 7 exact high-degree masters $\mathcal{M}^+_{(2,2,2)},\ldots,\mathcal{M}^+_{(8,0,0)}$ primed from \texttt{CubeT6Masters.wl} continue to anchor the leading parity-even contributions in closed form; PSLQ reconstruction of the remaining $\alpha_i,\beta_j$ in terms of $\sqrt{3}$, $\pi$, $\log$, $\tan^{-1}$, and $\sinh^{-1}$ atoms is the next analytical stage.


\subsection{High-precision scalar values}
\label{sec:results-high-precision}

{\sloppy
The values below are produced by \texttt{Mathematica/\allowbreak CubeT6OvernightHighPrec.wl} at \texttt{WorkingPrecision}\allowbreak$\to 30$ and stored verbatim in \texttt{Mathematica/\allowbreak CubeT6ScalarValues\_\allowbreak HighPrec.wl}.  Each entry carries $\ge 28$ correct decimal digits and is suitable as input to PSLQ-style integer-relation searches against an atom basis $\{1,\sqrt{3},\pi,\log 2,\log(1+\sqrt{3}),\log(\sqrt{3}-1),\sinh^{-1}1,\tan^{-1}(\cdot),\ldots\}$ in the spirit of Tables~\ref{tab:tier1-masters}--\ref{tab:tier2-masters}.\par}

\begin{table}[h]
\centering\small
\begin{tabular}{lll}
\toprule
Symbol & Value (28 digits) & Source orbit \\
\midrule
$\alpha_1$ & $\phantom{-}17.1263961354644384869032245660$ & $(10,10),(11,11)$ \\
$\alpha_2$ & $\phantom{-}14.6986795239648355731322989855$ & $(10,11),(11,12)$ \\
$\alpha_3$ & $\phantom{-1}3.9886969699081621215723702580$ & $(13,13),(14,14)$ \\
\midrule
$\beta_1$ & $\phantom{-1}7.0863777548574787932169236859$ & $(10,10)$ \\
$\beta_2$ & $\phantom{-1}5.0200091903034773113689403067$ & $(11,11)$ \\
$\beta_3$ & $\phantom{-1}5.4449549083142088200682904915$ & $(10,11)$ \\
$\beta_4$ & $\phantom{-1}3.8087697073364094814150175581$ & $(11,12)$ \\
$\beta_5$ & $\phantom{-1}1.6440754348395546266255471889$ & $(13,13)$ \\
$\beta_6$ & $\phantom{-1}1.1723107675343058603685866457$ & $(14,14)$ \\
$\beta_8$ & $\phantom{-1}0.6541891913383859981036052975$ & $(10,17),(11,16),(14,26)$ \\
$\beta_9$ & $-0.3895785442629410683383767348$ & $(13,20),(14,19)$ \\
$\beta_{10}$ & $-0.2164145776711156168035387731$ & $(10,21),(11,21)$ \\
\bottomrule
\end{tabular}
\caption{High-precision values of the $12$ independent scalars determining the $18\times 18$ quad--quad block of $\mathcal{B}_{27}^{\text{quad}}$, computed at \texttt{WorkingPrecision}$\to 30$.  The previously reported $\beta_7\approx 0.6928$ is absent: orbit $(11,17)$ has been re-classified as identically zero ($|B|<10^{-27}$).  Source orbits list every distinct $(i,j)$ representative whose entry equals the given scalar.}
\label{tab:high-prec-scalars}
\end{table}

\input{cube_galerkin27_closedforms.tex}

\paragraph{Per-orbit timing.}
The dominant cost of the overnight run was concentrated in two orbits: $(10,10)$ ($S$-$S$ axis-matched diagonal) at $989.6$ s, and $(10,11)$ ($S$-$S$ off-diagonal one-matched) at $140.3$ s.  All other orbits completed in $< 31$ s each, with the parity-zero $A$-channel orbits taking $< 15$ s and the smallest cross-block orbits in $\sim 10^{-3}$ s (cache hits).  Total wall time was $1217.07$ s, dominated by the single $(10,10)$ entry.

\paragraph{Verification.}
{\sloppy
Cubic-symmetry consistency for the 13 surviving orbits was checked at \texttt{WorkingPrecision}\allowbreak$\to 30$ and matches to $< 10^{-27}$ in all six independent identities.  The within-orbit spreads collapse from the $\sim 10^{-15}$ floating-point noise of the \texttt{MachinePrecision} run to $\sim 10^{-28}$, consistent with all displayed digits being algebraically exact.  Independent reproduction of $\alpha_1$ through $\alpha_3$ via the closed-form Tier-1 master integrals $\mathcal{M}_{p,q,r}$ (Table~\ref{tab:tier1-masters}) and the linear-combination assembly recipe of \eqref{eq:B27quad-split} agrees with the high-precision values to all $28$ digits, completing the cross-check.\par}


\subsection{Elastic stiffness bilinear form
\texorpdfstring{$\mathcal{B}_{\mathrm{el}}$}{B\_el} on
\texorpdfstring{$T_{27}$}{T\_27}}
\label{sec:results-Bel27}

Where the body bilinear form $\mathcal{B}_{\mathrm{body}}$ of Sections~\ref{sec:results-body9}--\ref{sec:results-body27quad} carries the (singular) Green-tensor double integral, the \emph{elastic} bilinear form $\mathcal{B}_{\mathrm{el}}$ is the ordinary strain-energy inner product of the isotropic stiffness contrast $\Delta c$ against the trial-space basis.  For the isotropic Kelvin form
\begin{equation}
  \Delta c_{ijkl}
  \;=\; \Delta\lambda\,\delta_{ij}\delta_{kl}
     \;+\; \Delta\mu\,\bigl(\delta_{ik}\delta_{jl} + \delta_{il}\delta_{jk}\bigr),
  \label{eq:deltaC-isotropic}
\end{equation}
the Galerkin matrix entries are the polynomial integrals
\begin{equation}
  \mathcal{B}_{\mathrm{el}}[v, w]
  \;=\; \int_{[-1,1]^3}\!\varepsilon(v):\Delta c:\varepsilon(w)\,dV
  \;=\; \int_{[-1,1]^3}\!\Bigl[\Delta\lambda\,\operatorname{tr}\varepsilon(v)\operatorname{tr}\varepsilon(w)
     \;+\; 2\Delta\mu\,\varepsilon(v){:}\varepsilon(w)\Bigr]\,dV,
  \label{eq:Bel-definition}
\end{equation}
with $\varepsilon(v)_{ij} = \tfrac12(\partial_i v_j + \partial_j v_i)$.  Because the $27$-dimensional trial space is polynomial, \eqref{eq:Bel-definition} is a closed-form symbolic integral and the full $27\times 27$ matrix $\mathcal{B}_{\mathrm{el}}$ is produced directly in \texttt{Mathematica/CubeT27Assemble.wl} Section~11 (no hand-coded block entries).

\paragraph{Structural result: the linear--quadratic cross block vanishes by parity.}
The decisive structural observation is that the $9\times 18$ block $\mathcal{B}_{\mathrm{el}}[\,1{:}9,\,10{:}27\,]$ connecting the linear subspace $T_9$ to the quadratic subspace $T_{27}^{\mathrm{quad}}$ is \emph{identically zero}:
\begin{equation}
  \mathcal{B}_{\mathrm{el}}\bigl[\,T_9,\,T_{27}^{\mathrm{quad}}\,\bigr]
  \;=\; 0.
  \label{eq:Bel-cross-zero}
\end{equation}
The mechanism is elementary.  The strain of any linear basis vector in $T_9$ is a \emph{constant} tensor, whereas the strain of any quadratic basis vector in $T_{27}^{\mathrm{quad}}$ is \emph{linear} in the coordinates $(r_1,r_2,r_3)$.  The integrand in \eqref{eq:Bel-definition} is therefore an odd polynomial on the symmetric box $[-1,1]^3$, and every component of every such cross entry integrates to $0$ term-by-term.  This cancellation is exact and algebraic, not numerical; the symbolic Mathematica run in \texttt{CubeT27Assemble.wl} confirms it by computing the $\sup$-norm of the full $9{\times}18$ block:
\begin{center}\texttt{|BelSym[[1..9,\,10..27]]|\ max\ =\ 0}.\end{center}

The immediate Galerkin-closure consequence is that $\mathcal{B}_{\mathrm{el}}$ is \emph{block-diagonal} with respect to the decomposition $T_{27} = T_9 \oplus T_{27}^{\mathrm{quad}}$:
\begin{equation}
  \mathcal{B}_{\mathrm{el}}
  \;=\;
  \begin{pmatrix}
    \mathcal{B}_{\mathrm{el}}^{\,T_9}   & 0 \\[2pt]
    0 & \mathcal{B}_{\mathrm{el}}^{\,\text{quad--quad}}
  \end{pmatrix}.
  \label{eq:Bel-block-diagonal}
\end{equation}
This stands in sharp contrast to the body form $\mathcal{B}_{\mathrm{body}}$, whose cross block is \emph{not} zero (see the Section~\ref{sec:results-body27quad} YY block and the $T_9{\times} T_{27}^{\mathrm{quad}}$ couplings that drive the body-channel Schur complement).

\paragraph{The $T_9$ local block.}
On the $9{\times}9$ linear sub-block, constants contribute no strain and the $3{\times}3$ displacement block is identically zero.  The remaining $6{\times}6$ strain sub-block (rows/cols $4\ldots 9$) reproduces the classical local Kelvin form $V\cdot\varepsilon{:}\Delta c{:}\varepsilon$ with $V=8$:
\begin{equation}
  \mathcal{B}_{\mathrm{el}}^{\,T_9}\bigl[\,4{:}9,\,4{:}9\,\bigr]
  \;=\; 8
  \begin{pmatrix}
    \Delta\lambda+2\Delta\mu & \Delta\lambda & \Delta\lambda & 0 & 0 & 0\\
    \Delta\lambda & \Delta\lambda+2\Delta\mu & \Delta\lambda & 0 & 0 & 0\\
    \Delta\lambda & \Delta\lambda & \Delta\lambda+2\Delta\mu & 0 & 0 & 0\\
    0 & 0 & 0 & \Delta\mu & 0 & 0\\
    0 & 0 & 0 & 0 & \Delta\mu & 0\\
    0 & 0 & 0 & 0 & 0 & \Delta\mu
  \end{pmatrix}.
  \label{eq:Bel-T9-strain}
\end{equation}
The $\mathrm{O}_h$-irreducible eigenvalues of \eqref{eq:Bel-T9-strain} are
\begin{align}
  \mathcal{B}_{\mathrm{el}}^{\,T_9}\Big|_{A_{1g}} \;&=\; 8\bigl(3\Delta\lambda + 2\Delta\mu\bigr)
  \;=\; 8\cdot 3K_{\Delta}, \label{eq:Bel-T9-A1g}\\
  \mathcal{B}_{\mathrm{el}}^{\,T_9}\Big|_{E_{g}}   \;&=\; 16\,\Delta\mu,                       \label{eq:Bel-T9-Eg}\\
  \mathcal{B}_{\mathrm{el}}^{\,T_9}\Big|_{T_{2g}}  \;&=\; \phantom{1}8\,\Delta\mu,             \label{eq:Bel-T9-T2g}
\end{align}
where $K_{\Delta} := \Delta\lambda + \tfrac23\Delta\mu$ is the bulk contrast.  The $T_9$ strain block carries exactly \emph{two} independent scalars $(\Delta\lambda,\Delta\mu)$ -- it is fully isotropic -- whereas the Path-A coefficients $(T_{1c},T_{2c},T_{3c})$ of the cube effective $T$-matrix carry \emph{three} independent scalars, the third one being the cubic-anisotropy piece associated with the static Eshelby integral $C$ (see the Path-A reference in \texttt{cubic\_scattering/\allowbreak effective\_contrasts.py}).  Equations \eqref{eq:Bel-T9-A1g}--\eqref{eq:Bel-T9-T2g} make it immediate that the missing cubic anisotropy cannot come from the stiffness channel in isolation.

\paragraph{The quad--quad $18{\times}18$ block.}
On the $18{\times}18$ sub-block $\mathcal{B}_{\mathrm{el}}^{\,\text{quad--quad}}$ with the \texttt{basisT27quad} ordering of Section~\ref{sec:results-body27quad}, the symbolic integration yields a matrix with $54$ non-zero entries out of $324$ (structurally sparse, diagonal-dominant) and symbolic rank $18$.  Per-direction the $6{\times}6$ diagonal splits into four distinct scalar entries, reflecting whether the quadratic monomial is of $S$ or $X$ type and whether its factor axes match the direction index $k$:
\begin{equation}
  \operatorname{diag}\!\bigl(\mathcal{B}_{\mathrm{el}}^{\,\text{quad--quad}}\bigr)_{\text{dir.}\,k}
  \;=\;
  \begin{cases}
    \tfrac{32}{3}\bigl(\Delta\lambda + 2\Delta\mu\bigr)
      & q = r_k^2 \quad\text{($S$, matched)},\\[6pt]
    \tfrac{32}{3}\,\Delta\mu
      & q = r_m^2,\ m\neq k \quad\text{($S$, unmatched)},\\[6pt]
    \tfrac{16}{3}\,\Delta\mu
      & q = r_m r_n,\ m,n\neq k \quad\text{($X$, no match)},\\[6pt]
    \tfrac{8}{3}\bigl(\Delta\lambda + 3\Delta\mu\bigr)
      & q = r_k r_m,\ m\neq k \quad\text{($X$, one match)}.
  \end{cases}
  \label{eq:Bel-quad-diagonal}
\end{equation}
The four per-direction diagonals \eqref{eq:Bel-quad-diagonal} are the analogues, on the quadratic sub-basis, of the three $T_9$ eigenvalues \eqref{eq:Bel-T9-A1g}--\eqref{eq:Bel-T9-T2g}.  Their orbit structure under $\mathrm{O}_h$ matches exactly the parity classification of Table~\ref{tab:tier6-orbits}.

\paragraph{Consequence for the Schur-complement reduction to $T_9$.}
Under the body-channel Galerkin eigenproblem $(\mathcal{M}_{27} - \varepsilon\,\mathcal{B})\,\mathbf{c} = 0$ the Schur-complemented effective operator on $T_9$ is
\begin{equation*}
  T_9^{\mathrm{eff}}
  \;=\; \bigl(\mathcal{M}_{11} - \mathcal{M}_{12}\mathcal{M}_{22}^{-1}\mathcal{M}_{21}\bigr)^{-1}
     \bigl(\mathcal{B}_{11} \;-\; \mathcal{M}_{12}\mathcal{M}_{22}^{-1}\mathcal{B}_{21} \;-\; \mathcal{B}_{12}\mathcal{M}_{22}^{-1}\mathcal{M}_{21}\bigr).
\end{equation*}
For the stiffness channel, $\mathcal{B} = \mathcal{B}_{\mathrm{el}}$ and \eqref{eq:Bel-cross-zero} forces $\mathcal{B}_{12} = \mathcal{B}_{21}^{\top} = 0$, so the middle and trailing terms vanish identically and the Schur-complemented stiffness operator reduces to the local block \eqref{eq:Bel-T9-strain}:
\begin{equation}
  T_9^{\mathrm{eff,\,el}}
  \;=\; \bigl(\mathcal{M}_{11} - \mathcal{M}_{12}\mathcal{M}_{22}^{-1}\mathcal{M}_{21}\bigr)^{-1}\,
     \mathcal{B}_{\mathrm{el}}^{\,T_9}\bigl[\,1{:}9,\,1{:}9\,\bigr].
  \label{eq:T9-eff-stiffness-schur}
\end{equation}
In particular, the entire $T_{27}^{\mathrm{quad}}$ sub-block of $\mathcal{B}_{\mathrm{el}}$ drops out of the $T_9$ reduction at leading order, and the cubic-anisotropy scalar $T_{3c}$ of the Path-A effective $T$-matrix \emph{cannot} originate in the elastic Galerkin channel on the quadratic sub-basis.  The parallel observation of Section~\ref{sec:results-body9} was that the body-channel $T_9^{\mathrm{eff}}$ acts as a pure scalar on the $6$-dim strain sector in the numerical $A$-polarization probe of Section~12 (Schur-complemented diagonals \texttt{const = -47.1737...}, \texttt{strain = 5.7647...}, off-diagonals exactly $0$ at machine precision).  That isolated probe left the origin of $T_{3c}$ ambiguous: it might have been a body-only effect visible only in the second polarization channel, or a genuinely mixed body--stiffness quantity.  The symbolic computation in Section~12b below resolves this ambiguity by keeping both polarization channels explicit, and shows that $T_{3c}$ arises from the $B$-polarization of the body operator alone (see \eqref{eq:T3c-signature}); the stiffness channel contributes only isotropic corrections.

\subsubsection*{Symbolic verification (Section 11b)}
\label{sec:results-bel-schur-11b}
An explicit symbolic computation of the stiffness Schur complement has been added as Section~11b of \texttt{CubeT27Assemble.wl} and run end-to-end at exact rational precision.  Partitioning the $27{\times}27$ matrices into $9\mid 18$ blocks
\[
  \mathcal{M}_{27} = \begin{pmatrix}\mathcal{M}_{11} & \mathcal{M}_{12}\\ \mathcal{M}_{21} & \mathcal{M}_{22}\end{pmatrix},
  \qquad
  \mathcal{B}_{\mathrm{el}} = \begin{pmatrix}\mathcal{B}_{11}^{\mathrm{el}} & \mathcal{B}_{12}^{\mathrm{el}}\\ \mathcal{B}_{21}^{\mathrm{el}} & \mathcal{B}_{22}^{\mathrm{el}}\end{pmatrix},
\]
and forming the full Schur complement
\[
  \mathcal{B}_{\mathrm{el}}^{\mathrm{Schur}}
  \;\equiv\; \mathcal{B}_{11}^{\mathrm{el}}
     - \mathcal{M}_{12}\mathcal{M}_{22}^{-1}\mathcal{B}_{21}^{\mathrm{el}}
     - \mathcal{B}_{12}^{\mathrm{el}}\mathcal{M}_{22}^{-1}\mathcal{M}_{21},
\]
Mathematica returns $\mathsf{Simplify}\bigl[\mathcal{B}_{\mathrm{el}}^{\mathrm{Schur}} - \mathcal{B}_{11}^{\mathrm{el}}\bigr] = 0$ identically (not just numerically), confirming \emph{algebraically} that the parity-vanishing result $\mathcal{B}_{12}^{\mathrm{el}} = (\mathcal{B}_{21}^{\mathrm{el}})^{\top} = 0$ makes the middle and trailing Schur terms disappear exactly.  The strain-sector (rows/cols $4\dots 9$) eigenvalues of $\mathcal{B}_{\mathrm{el}}^{\mathrm{Schur}}$ are then the pure isotropic Kelvin form,
\begin{equation}
  \lambda_{A_{1g}} \;=\; 8\,(3\,\Delta\lambda + 2\,\Delta\mu),
  \qquad
  \lambda_{E_{g}} \;=\; 16\,\Delta\mu,
  \qquad
  \lambda_{T_{2g}} \;=\; 8\,\Delta\mu,
  \label{eq:Bel-Schur-eigs}
\end{equation}
matching the direct diagonal extraction \eqref{eq:Bel-T9-A1g}--\eqref{eq:Bel-T9-T2g} of Section~11.  Post-multiplying by the body-channel mass Schur inverse to form the full stiffness-channel effective operator on $T_9$,
\[
  T_9^{\mathrm{eff,\,el}}
  \;=\; \bigl(\mathcal{M}_{11} - \mathcal{M}_{12}\mathcal{M}_{22}^{-1}\mathcal{M}_{21}\bigr)^{-1}\,\mathcal{B}_{\mathrm{el}}^{\mathrm{Schur}},
\]
the symbolic simplification gives strain-sector irrep eigenvalues
\begin{equation}
  \lambda_{A_{1g}}\bigl(T_9^{\mathrm{eff,\,el}}\bigr) = 9\,\Delta\lambda + 6\,\Delta\mu,
  \qquad
  \lambda_{E_{g}}\bigl(T_9^{\mathrm{eff,\,el}}\bigr) = 6\,\Delta\mu,
  \qquad
  \lambda_{T_{2g}}\bigl(T_9^{\mathrm{eff,\,el}}\bigr) = 6\,\Delta\mu.
  \label{eq:T9-eff-el-eigs}
\end{equation}
The crucial observation is the equality
\begin{equation}
  \lambda_{E_g}\bigl(T_9^{\mathrm{eff,\,el}}\bigr)
  \;=\; \lambda_{T_{2g}}\bigl(T_9^{\mathrm{eff,\,el}}\bigr)
  \;=\; 6\,\Delta\mu,
  \label{eq:Eg-T2g-degeneracy}
\end{equation}
which holds \emph{symbolically} in $\Delta\lambda, \Delta\mu$ with no residual.  Cubic anisotropy, by definition, requires $\lambda_{E_g} \neq \lambda_{T_{2g}}$ on the five-dimensional deviatoric strain sector (the $E_g \oplus T_{2g}$ splitting); the fact that they are \emph{exactly equal} in $T_9^{\mathrm{eff,\,el}}$ therefore rules out any cubic-anisotropic component from the stiffness channel in isolation.  Equivalently, $T_9^{\mathrm{eff,\,el}}$ is strictly proportional to the isotropic Kelvin contraction, and its Path-A decomposition acquires no contribution to $T_{3c}$.  The complementary question -- whether the body channel by itself already produces a nonzero $T_{3c}$, or whether cross-coupling to $\mathcal{B}_{\mathrm{el}}$ is required -- is answered symbolically in Section~12b below.

\subsubsection*{Origin of $T_{3c}$: combined body+stiffness Schur (Section 12b)}
\label{sec:results-body-stiffness-12b}

With the stiffness channel definitively isotropic \eqref{eq:Eg-T2g-degeneracy}, the origin of the cube's $T_{3c}$ cubic-anisotropy scalar must lie in the body channel.  Section~12b of \texttt{CubeT27Assemble.wl} completes this identification by symbolically constructing the combined body+stiffness Schur-complemented effective operator
\begin{equation}
  T_9^{\mathrm{eff,\,tot}}
  \;=\; \bigl(\mathcal{M}_{11}
     - \mathcal{M}_{12}\mathcal{M}_{22}^{-1}\mathcal{M}_{21}\bigr)^{-1}
     \Bigl[\eta_{\mathrm{body}}\,\mathcal{B}_{\mathrm{body}}^{\mathrm{Schur}}
       + \mathcal{B}_{\mathrm{el}}^{\mathrm{Schur}}\Bigr],
  \label{eq:T9-eff-total-schur}
\end{equation}
where $\eta_{\mathrm{body}}$ is a free body-coupling coefficient (physically $\propto \omega^{2}\Delta\rho$) and $\mathcal{B}_{\mathrm{body}}^{\mathrm{Schur}}$ is the $9{\times}9$ Schur complement of the full $27{\times}27$ body bilinear form.  The key technical step is to keep $\mathcal{B}_{\mathrm{body}}^{\mathrm{Schur}}$ symbolic in the two independent polarization channels $(A_{\mathrm{elas}}, B_{\mathrm{elas}})$ of the underlying elastic Green's tensor:
\begin{itemize}
  \item The $A$-channel (coefficient $A_{\mathrm{elas}}$) collects the isotropic $\delta_{ij}$ piece of the Green's tensor, weighted by the $K_{1}$ kernel (weak $1/|r|$ integrals over the Mp masters of Tier-6).
  \item The $B$-channel (coefficient $B_{\mathrm{elas}}$) collects the tensorial anisotropic piece, i.e.\ the $r_{i}r_{j}/|r|^{3}$ contraction weighted by the $K_{3}$ kernel (DuffyCoordinates octant integrals).
\end{itemize}
These correspond to the two independent building blocks $(\mathsf{aPiece},\mathsf{bPiece})$ returned by \texttt{bbBodyAB} in \texttt{CubeT6Block.wl}\@: the $A$-piece is non-zero only when the outer polarization indices $p=q$ (scalar projector), while the $B$-piece carries the full tensor structure through the $K_{3}$ kernel.

Substituting the Pell-simplified atomic scalars (\texttt{atomRules}) while holding $(A_{\mathrm{elas}}, B_{\mathrm{elas}})$ symbolic, Mathematica returns for $\mathcal{B}_{\mathrm{body}}^{\mathrm{Schur}}$ restricted to the strain sector (rows/cols $4\dots 9$) the irrep eigenvalues
\begin{equation}
\begin{aligned}
  \lambda_{A_{1g}}\bigl(\mathcal{B}_{\mathrm{body}}^{\mathrm{Schur}}\bigr) &= 15.3725\,A_{\mathrm{elas}} + 1.0059\,B_{\mathrm{elas}}, \\
  \lambda_{E_{g}}\bigl(\mathcal{B}_{\mathrm{body}}^{\mathrm{Schur}}\bigr) &= 15.3725\,A_{\mathrm{elas}} + 5.4217\,B_{\mathrm{elas}}, \\
  \lambda_{T_{2g}}\bigl(\mathcal{B}_{\mathrm{body}}^{\mathrm{Schur}}\bigr) &= \phantom{0}7.6863\,A_{\mathrm{elas}} + 2.1197\,B_{\mathrm{elas}},
\end{aligned}
\label{eq:Bbody-Schur-eigs}
\end{equation}
all to $\sim\!16$-digit MachinePrecision.  Both channels exhibit $\lambda_{E_g} \neq \lambda_{T_{2g}}$ at this raw Schur-complement stage, so a naive cubic-anisotropy test on $\mathcal{B}_{\mathrm{body}}^{\mathrm{Schur}}$ alone would (misleadingly) flag both channels.  However, the physically relevant operator is the full effective $T_9^{\mathrm{eff,\,tot}}$ obtained after pre-multiplying by the mass Schur inverse.  The symbolic simplification gives
\begin{equation}
\begin{aligned}
  \lambda_{A_{1g}}\bigl(T_9^{\mathrm{eff,\,tot}}\bigr) &= 9\,\Delta\lambda + 6\,\Delta\mu + \bigl(5.7647\,A_{\mathrm{elas}} + 0.3772\,B_{\mathrm{elas}}\bigr)\,\eta_{\mathrm{body}}, \\
  \lambda_{E_{g}}\bigl(T_9^{\mathrm{eff,\,tot}}\bigr) &= \phantom{9\,\Delta\lambda + {}}6\,\Delta\mu + \bigl(5.7647\,A_{\mathrm{elas}} + 2.0331\,B_{\mathrm{elas}}\bigr)\,\eta_{\mathrm{body}}, \\
  \lambda_{T_{2g}}\bigl(T_9^{\mathrm{eff,\,tot}}\bigr) &= \phantom{9\,\Delta\lambda + {}}6\,\Delta\mu + \bigl(5.7647\,A_{\mathrm{elas}} + 1.5898\,B_{\mathrm{elas}}\bigr)\,\eta_{\mathrm{body}}.
\end{aligned}
\label{eq:T9-eff-total-eigs}
\end{equation}
The $A$-channel coefficients on $\eta_{\mathrm{body}}$ are now \emph{all equal} to $5.7647$: the anisotropy visible in \eqref{eq:Bbody-Schur-eigs} in the $A$-polarization cancels exactly against the equivariant mass-Schur inverse, reproducing isotropy.  In the $B$-channel the coefficients split as $(0.3772,\,2.0331,\,1.5898)$, yielding the cubic-anisotropy signature
\begin{equation}
  \lambda_{E_{g}}\bigl(T_9^{\mathrm{eff,\,tot}}\bigr) - \lambda_{T_{2g}}\bigl(T_9^{\mathrm{eff,\,tot}}\bigr)
  \;=\; 0.4433\,B_{\mathrm{elas}}\,\eta_{\mathrm{body}}
  \;\neq\; 0,
  \label{eq:T3c-signature}
\end{equation}
which is \emph{independent} of both the stiffness contrasts $(\Delta\lambda, \Delta\mu)$ and the $A$-polarization of the body operator.

Equation \eqref{eq:T3c-signature} is the positive identification of the cube's $T_{3c}$ anisotropy at leading order in the $T_{27}$ Galerkin closure.  Three consequences deserve emphasis:
\begin{enumerate}
  \item The ``stiffness-only'' probe $\eta_{\mathrm{body}} = 0$ gives $\lambda_{E_g} - \lambda_{T_{2g}} = 0$ (Section~11b result), confirming no cubic-anisotropy component from $\mathcal{B}_{\mathrm{el}}$ alone.
  \item The ``$A$-only body'' probe $B_{\mathrm{elas}} = 0$ also gives $\lambda_{E_g} - \lambda_{T_{2g}} = 0$: the trace/isotropic piece $\delta_{ij}$ of the elastic Green's tensor (K$_{1}$ kernel) is insufficient, even when fully Galerkin-coupled through $T_{27}^{\mathrm{quad}}$.
  \item The $B$-channel $r_{i}r_{j}/|r|^{3}$ tensor piece of the elastic Green's tensor (K$_{3}$ kernel) is the sole source of the leading-order $T_{3c}$ signature in the tier-$27$ closure.
\end{enumerate}

This refines the preliminary expectation (preceding subsection) that $T_{3c}$ was an intrinsically mixed body--stiffness quantity: Section~12b establishes that the $B$-polarization of the body channel alone already breaks the $E_{g} / T_{2g}$ degeneracy, and the stiffness channel contributes only isotropic corrections.  The numerical coefficient $0.4433\,B_{\mathrm{elas}}\,\eta_{\mathrm{body}}$ is a closed form awaiting PSLQ/Pell simplification against the same basis of elementary constants used for the Tier-6 scalars $\{\alpha_{i}\}_{i=1}^{3}$ and $\{\beta_{j}\}_{j=1}^{9}$ of \texttt{CubeT6ScalarValues\_HighPrec\_Pell.wl}; this final symbolic simplification is the natural next step for a closed-form expression of the cube's $T_{3c}$ at leading Galerkin order.

